\section{Conclusiones y trabajo futuro}

Este trabajo presentó una comparación empírica de cuatro Small Language
Models open-source (0,36B a 1,7B parámetros) en la tarea de interpretar
comandos de domótica en español rioplatense, usando un dataset propio de
32 comandos anotados. Los resultados muestran que la coincidencia exacta
mejora de forma consistente con el tamaño del modelo (18,8\% a 59,4\%),
pero con una latencia que crece más rápido que la exactitud, y que los
cuatro modelos comparten una debilidad sistemática en la interpretación de
comandos con ajustes relativos o cualitativos sin valor numérico explícito,
tendiendo a alucinar valores y unidades inexistentes en el comando
original — un patrón que no se corrige con más escala dentro del rango
evaluado.

Como trabajo futuro se plantea: (1) ampliar el dataset con más comandos por
categoría y con voz real transcripta por ASR; (2) evaluar el efecto de la
cuantización agresiva (4-bit, 8-bit) sobre el trade-off exactitud/latencia
en el mismo hardware; (3) evaluar estrategias de mitigación del sesgo de
alucinación de valor/unidad, como post-procesamiento basado en reglas o
fine-tuning de dominio con ejemplos negativos explícitos; y (4) extender la
evaluación a comandos multi-turno con contexto de diálogo.
